% Nouveau EGA V (cote 25): archive cote 25, site page 119 (PDF page 120).
% typescript. Draft compared: none.
% AI-transcribed by Claude Opus 5.5 (claude-opus-5-5, adaptive thinking, effort high), Message Batches API, batch ega25b (msgbatch_01A7n5YuXubUPNK1zsepi8RG), 2026-11-13; not yet checked by eye.
\scanpage{25/119}
\margin{\unsure{2/2}}
\margin{(trait vertical au crayon dans la marge gauche, le long du premier paragraphe)}
Je me répens d'avoir donné la proposition \deleted{\unsure{3}} sous une forme éculée et
pisseuse, à mi chemin entre  les hypothèses classiques et les hypothèses
naturelles, et sans donner une récip\unsure{r}oque. Aussi je propose l'énoncé:
\uline{Proposition\unsure{\hand{3}}3. \deleted{revue et corrigée}}.  Soit $X \to S$ plat loc de présentation
finie,  $Q=\check{P} = P(\check{E})$  un fibré \xout{\unsure{vectoriel sur}} projectif sur $S$ défini par
un Module loc  libre de type fini $F = \check{E}$ ,  considérons l'ensemble
$\emptyset$  des systêmes linéaires \inserted{$D$ \deleted{\unsure{E}}} de diviseurs \xout{\unsure{sur}} sur $X$, paramétrés par $Q$, tels
que l'ensemble  $Z$ des points fixes de $D$  satisfasse la propriété: $z \in Z$
implique prof$_z \underline{O}_{X_s} \geqslant 2$  (où  $s$ \xout{\unsure{$= g(z)$)}} est l'image de $z$ dans $S$). \xout{\unsure{Alors}}
\xout{\unsure{$D \in \emptyset$ associons le morphisme correspondant $f$ $X$-$Z$ $P$ considérée}}
\xout{\unsure{considérée comme}}  Soit $\emptyset$ l'ensemble des  \deleted{``applications rationelles}\inserted{pseudo-morphismes} rel à
$S$ \deleted{\unsure{*}}\inserted{$f$} de $X$ dans  $P$, telles que le\xout{\illegible{1}} domaine de définition \inserted{$U=$}\over{dom}{\unsure{U}}(f)
\xout{\unsure{f satisf}} satisfasse la condition : $z \in X$-$U$\deleted{(f)} $\Longrightarrow$ prof$_z \underline{O}_{X_s} \geqslant 2$\inserted{$\surd$, et $f_U = f|U : U \to \check{P}$ satisfait la condition de non dégénérescence envisagée dans 3.1.}. Considé\xout{rons}
rons l'application naturelle  $D \rightsquigarrow f_D$  de $\emptyset$  dans $\emptyset$ .  Alors :

a)  Cette application est injective,  et  pour $D \in \emptyset$ , l'ensemble des
points fixes $Z$  n'est autre que l'ensemble  complémentaire du domaine de
définition \xout{\unsure{Dif}} $U$\deleted{(f)} de $f$\hand{$_D$}.

b)  \xout{\unsure{Pour qu'une f $\emptyset$ appartienne à provienne d'une $D$ $\emptyset$, il}} Soit
$f \in \emptyset$ , et soit $U$ un ouvert sur lequel $f$ soit défini, et tel que $z \in X$-$U$
implique prof $\underline{O}_{X_s,z} \geqslant 2$ \xout{\unsure{et existe}} (par exemple $U = U(f)$ domaine de
définition de $f$). Pour que $f$ provienne d'une $D \in \emptyset$ , il f et suffit  que\hand{\unsure{$\surd$}}
\deleted{\unsure{$f_U$ satisfasse la condition énoncée dans 3.1. et que}}
posant $\underline{L}_U = f_U^*(\underline{O}_P(1))$  (où  $f_U$: $U \to P$  est le morphisme induit par $f$),
\deleted{\unsure{$f_U$ satisfasse la condition énoncée dans 3.1 et}}
\xout{\unsure{que}} le Module  $i_*(\underline{L}_U)$  sur $X$ soit inversible (où $i$:$U \to X$ est l'immer-%
sion canonique).

On remarquera que lorsque  les fibres de $X$ sur $S$ satisfont à $(S_2)$,
par exemple sont normales, voire géométriquement  normales, alors la
condition de profondeur  envisagée sur  une partie fermée $Z$ de $X$ dans \deleted{\illegible{1}}
la proposition  signifie simplement que pour tout $s \in S$,

\mspage{119}
